\scanpage{24}
\mspage{7}
de $\dot X\subset\widetilde X$ est déterminé par la connaissance de
\[
\widetilde X'_1\cap\dot X=\dot X'_1.
\]

\uline{Corollaire} Les parties \inserted{fermées} $X'$ de $X$ compatibles avec
le germe de structure cônique, sont en
corr. $1$-$1$ avec les couples $(\widetilde X'\subset\widetilde X,Y'\subset Y)$
de parties fermées de $\widetilde X$ et de $Y$, telles que

\unsure{a)|1°)} $\widetilde X'$ soit compatible avec le \inserted{germe de} structure cylindrique
\inserted{de $\widetilde X,\dot X$}

\unsure{b)|2°)} $f:\dot X\to Y$ applique \deleted{$\unsure{\widetilde X'|\widetilde X}$}
$\dot X'\overset{\text{déf}}{=}\widetilde X'\cap\dot X$ dans $Y'$.

On aura
% P07: the upper brace begins beneath the right half of Im's m;
% it is represented as attaching to the whole image expression.
\[
X'=\underbrace{\operatorname{Im}(\widetilde X'\ \text{par}\ \widetilde X\to X)}
_{\overline{X'\setminus\underbrace{Y\cap X'}_{Y'}}}\cup Y'
\]
Pour que $X'$ soit \deleted{\uline{strict\unsure{.}} compatible avec}
\inserted{cônique \uline{strict}},
% P09: the cancelled raised word attaches after the old article, before structure.
\deleted{\unsure{la|le} \inserted{\unsure{purement|germe de}} structure cônique}
il f et. s. que $Y'=\operatorname{Im}(\dot X'\to Y)$

Soit maintenant $X_0\subset X$ une partie
bordante de $X$, et $\mathcal E$ une \inserted{germe de} structure cylindrique
de $X_0$ dans $X$. On dit que $\mathcal E$ est compatible
avec le germe de structure conique $\mathcal T$
de $X$ autour de $Y$, si $\exists$ représentants
$E\inserted{\simeq X_0\times\mathbf I}$ voisinage cylindrique de $X_0$, \unsure{\&|et} $T$ voisinage
cônique de $Y$ (avec \inserted{\unsure{sa|la}} structure cônique), tels que

\margin{au lieu de ``compatible'' il est plus parlant de
dire ``\uline{transversal}''}

a) $T\cap\unsure{\overline E|E}$, $\dot T\cap E$, et $Y\cap\unsure{\overline E|E}$
\deleted{\unsure{compat.}} \inserted{soient} des parties cylindriques de
\deleted{$\unsure{\overline E|E}\cap\unsure{\overline T|T}$ soit stable \unsure{par|pour} les deux types d'\unsure{homothéties}
$\unsure{\overline x|x}\mapsto\unsure{{\lambda x}_{X_0}|\lambda_{X_0}x}$ et
$t\unsure{\longrightarrow|\mapsto}\unsure{\lambda_Y\times t|\lambda\times_Y t}$
\unsure{qui|que} \unsure{agissent} \unsure{de part et d'autre}}
% P18: the second map's Y is lowered between lambda and the crossing sign;
% attachment to lambda is preferred but not settled. The final t is crossed.
\deleted{$E$ \unsure{i.e.}}

\margin{donc stable par $E$-homothéties par $\lambda\in[0,1]$}

% P20: the entire diagram is cancelled by nine long diagonal strokes.
% The leftmost source stroke also crosses the separately transcribed E i.e.
% Normal diagram spacing is used; no extra arrowheads are supplied to the cups.
\diagram{
\begin{tikzcd}[ampersand replacement=\&,
 column sep=large,row sep=large,
 execute at end picture={
  \draw[line width=.6pt] ([xshift=-2.4cm,yshift=.5cm]pA.north west) -- ([xshift=-1.1cm,yshift=-.4cm]pF.south west);
  \draw[line width=.6pt] ([xshift=-1.6cm,yshift=.8cm]pA.north west) -- ([xshift=-.2cm,yshift=-.3cm]pF.south west);
  \draw[line width=.6pt] ([xshift=-.8cm,yshift=.8cm]pA.north west) -- ([xshift=.6cm,yshift=-.4cm]pF.south west);
  \draw[line width=.6pt] ([xshift=-.3cm,yshift=.8cm]pA.north west) -- ([xshift=1.0cm,yshift=-.1cm]pF.south west);
  \draw[line width=.6pt] ([xshift=.2cm,yshift=.7cm]pA.north east) -- ([xshift=1.2cm,yshift=-.5cm]pF.south east);
  \draw[line width=.4pt] ([xshift=-.6cm,yshift=.8cm]pB.north west) -- ([xshift=.6cm,yshift=-.1cm]pG.south west);
  \draw[line width=.6pt] ([xshift=.1cm,yshift=.6cm]pB.north) -- ([xshift=1.5cm,yshift=-.5cm]pG.south);
  \draw[line width=.6pt] ([xshift=-.6cm,yshift=1.6cm]pE.north west) -- ([xshift=.5cm,yshift=-.2cm]pH.south west);
  \draw[line width=.6pt] ([xshift=.2cm,yshift=1.6cm]pE.north) -- ([xshift=1.4cm,yshift=-.4cm]pH.south);
 }
]
|[alias=pA]| E \arrow[r,phantom,"\simeq"] \&
|[alias=pB]| X_0\times I \\
|[alias=pC]| \dot T\cap\unsure{E|\overline E}
 \arrow[u,phantom,"\cup" description] \arrow[r,phantom,"\simeq"] \&
|[alias=pD]| T_0\times I \arrow[u,phantom,"\cup" description] \&
|[alias=pE]| T_0=X_0\cap T \\
|[alias=pF]| Y\cap\unsure{E|\overline E} \arrow[r,phantom,"\simeq"] \&
|[alias=pG]| Y_0\times I \arrow[u,phantom,"\cup" description] \&
|[alias=pH]| Y_0=\dot X_0\cap Y
\end{tikzcd}}
% P20: a separately crossed small group lies to the left of the lower diagram.
\deleted{$\unsure{E_0I|E\cap T}$}
\margin{23}
